\section{Discussion}
\label{sec:discussion}
The findings are, in essence, separating two costs that are amassing measurements combine; the consumer's remote free and the producer's
In mimalloc at 64~B the free tail is 250~ns and the allocation tail 1{,}904~ns, broadly speaking; allocation fail in separate places, broadly speaking.
In jemalloc the order is reversing, with a 1{,}040~ns allocation tail and a 3{,}232~ns free tail, at least empirically. A benchmark that reports
only a combined figure would not uncover which thread to tune, so also, the data do not establish the process behind either
split. They are, in this setting, clear with a remote-free path that is cheap in one design and a shared arrangement that is
claimed in the other, but no trial here separates that architecture.

Across the 16~B, the size dependence is, in essence, the most concrete finding; in turn, also, no person allocator ranking is holding to 256~KB range, all things considered.
For tcmalloc the large-size behaviour is a state the allocator is in for most iterations, not a homogeneous cost, in many respects. A larger
The transfer-batch cap is tested by a source-level trial; in turn, in a ten-run ablation, per-process thread-cache budget did not change the median, in real-world scenarios.
Raising the cap at 16~KB from four to 64 objects removes most of the slow mode (median allocation P99.9 from 737~$\mu$s to
74~$\mu$s). The same kind of change at 256~KB, from two to four objects, does not, so the plateau above 16~KB has a
differed cause, which the present examination did not spot, and the cap also trades small-object tails for large-object tails,
as the source comment says.

An supposition in current allocator design is changed by the coherence result, which is why the cost step for atomic writes on this ARM64 architecture
lies at a 64-byte stride. Padding from 64 to 128~B changes the measured cost by at most 0.8~ns and 1.7 cycles at either
thread count. The cycle counts disclose the cost is a processor stall, since the instruction count is flat, in actual application.
The micro-architectural penalty is, at least in law, directly quantified by hardware latency and cycle metrics, in real-world scenarios.

The release is, in effect, ensuing actionably guide for allocators managing virtual memory on modern UNIX kernels, in real-world scenarios. Without memory pressure,
\texttt{MADV\_DONTNEED} and standard \texttt{MADV\_FREE} do not change the size the framework reports. An allocator that rests
on them to bound size will see no reduction until pressure builds, which is why also, the put forward inquiry does not reproduce that
condition. \texttt{MADV\_FREE\_REUSABLE} does change size promptly. The access tails under the release were
indistinguishable from the no-release control, so the probe is finding no measurable interference from release on
concurrent is reading at this granularity. The reuse tails varied between repeat measurements. Also, No policy given a uniform reuse
advantage.

A second, separate issue is had by the system allocator, which is why still, its size under the KV-churn model was bimodal across the runs,
with half the runs near 10--14~MB and half between 40 and 88~MB; while the other four allocators reproduced their sizes
across repeat measurements. Any comparison that is applying a person measurement for the system allocator's size is therefore fragile, at least for now.
The cause is not spotted.
